\Article{365}
Sufficit nobis hic possibilitatem solutionum harum monstravisse. Multa artificia, quibus hae operationes sublevari possunt, praeterimus brevitatis gratia. At consequentias quasdam pergraves praetermittere non possumus.

Per praecedentia demonstratum est, omnes aequationes auxiliares pro solutione aequationis $x^\nu=1$, si in congruentias convertantur, habere radices possibiles, quando periodus
\SourceLeaf{244}
\[
x+x^p+x^{pp}+\dotstwo+x^{p^{m-1}}
\]
nondum est disiuncta. Subsistamus in casu, ubi $\nu$ est numerus primus; erit $m$ divisor ipsius $\nu-1$. Hic itaque congruentiae auxiliares, si numerus periodorum, quae per illas inveniuntur, est pars aliquota numeri $\frac{\nu-1}m$, habent radices reales. Si itaque $\frac{\nu-1}m$ est par i. e. si $m$ est divisor numeri $\frac{\nu-1}2$ seu si $p^{\frac{\nu-1}2}\equiv1\;(\mathrm{mod.}\,\nu)$ seu si $p$ est residuum quadraticum numeri primi $\nu$, aequatio quadratica, per quam radices in duas periodos dividuntur, habebit radices reales secundum modulum $p$. At in Cap. \textsc{vi} monstravimus, hanc aequationem posito $\nu=4n\pm1$ semper esse $xx+x\mp n=0$. Quare habetur insigne

\textsc{Theorema.} \emph{Si numerus primus $p$ est residuum quadraticum numeri primi $4n\pm1$, congruentia}
\[
xx+x\mp n\equiv0(\mathrm{mod.}\,p)
\]
\emph{habebit radices reales, adeoque etiam congruentia}
\[
4xx+4x\mp4n\equiv0\quad\text{seu}\quad(2x+1)^2\mp\nu\equiv0
\]
\emph{i. e. $\pm\nu$ erit residuum quadraticum numeri $p$.}

\Article{366}
Haec igitur est tertia theorematis fundamentalis Capitis \textsc{iv} completa demonstratio, eo magis attentione digna, quod principia, e quibus est petita, ab iis quibus ad priores usi sumus, prorsus sunt diversa.\ .At ex eodem hoc fonte, sed via opposita quartam deducamus. Scilicet sit $\nu$ numerus primus formae $4n\pm1$, $p$ alius primus quicunque, sitque $\pm\nu$ residuum quadraticum numeri primi $p$, demonstrabimus, $p$ fore residuum quadraticum numeri $\nu$.

Sit $p^m$ minima potestas numeri $p$, quae sit $\equiv1(\mathrm{mod.}\,\nu)$. Divisores elementares functionis $\frac{x^\nu-1}{x-1}$ secundum $p$ habebunt $m$ dimensiones, quare omnium numerus erit $=\frac{\nu-1}m$. Iam quoniam $\pm\nu Rp$, congruentia
\[
xx+x\mp n\equiv0(\mathrm{mod.}\,p)
\]
erit resolubilis; sint radices $A,A'$. Distribuantur functiones $x,xx,\dotstwo x^{\nu-1}$ in binas classes per $\xi,\xi'$ designandas, erit
\SourceLeaf{245}
\[
\begin{aligned}
\xi+\xi'&\equiv A+A'+(1+x+xx+\dotstwo+x^{\nu-1})\\
\xi\xi'&\equiv AA'+\lambda(1+x+xx+\dotstwo+x^{\nu-1})
\end{aligned}
\]
quare
\[
(z-\xi)(z-\xi')-(z-A)(z-A')
\]
per quemvis divisorem elementarem functionis $\frac{x^\nu-1}{x-1}$ erit divisibilis. Hinc autem quivis horum divisorum elementarium aut $\xi-A$ et $\xi'-A'$, aut $\xi-A'$ et $\xi'-A$ dividet. Hinc patet (quoniam $A$ non $\equiv A'$), si pro $x$ ponatur $x^p$, $\xi$ et $\xi'$ non immutari. Si enim $\xi$ in $\xi'$ et vice versa transiret, $\xi-A$ et $\xi-A'$ per eandem functionem primam dividerentur. Q. E. A. Hinc denique sequitur, $\frac{\nu-1}2$ per $m$ dividi seu $p^{\frac{\nu-1}2}-1$ per $\nu$. Quare $p$ erit residuum quadraticum ipsius $\nu$. Q. E. D.

Facile autem est omnes theorematis fundamentalis casus ex utroque theoremate derivare.

\Article{367}
Quamvis ad casum, ubi $\nu$ est numerus primus, hic nos restrinxerimus, tamen etiam, si $\nu$ sit compositus, theoremata analoga haud magno negotio determinari possunt, quod fusius exponere brevitatis gratia nunc non licet.

Manifestum est, similes observationes etiam de maiori periodorum multitudine formari posse. Ita si $\frac{\nu-1}m$ per $3$ dividitur i. e. si $p$ est residuum cubicum numeri primi $\nu$, aequatio, per quam radices aequationis $x^\nu=1$ in tres periodos distribuuntur quamque in Cap. \textsc{vi} a priori determinandam docuimus, solubilis erit secundum modulum $p$ et vice versa. Ita ex. gr. congruentia $x^3+xx-2x-1\equiv0$ secundum modulum primum quemcunque, qui est formae $7n\pm1$, resolvi potest, si vero aliam formam habeat, non poterit.

Non difficile nobis foret hoc Caput multis aliis observationibus locupletare, nisi limites, intra quos restringi oportet, vetarent. Iis qui ulterius progredi amant, haec principia viam saltem addigitare poterunt.
